% ============================================================
% Section 96: 抛物线色散金属的费米速度与电导率比较
% ============================================================
\section{固体理论：抛物线色散金属的费米速度与电导率比较}
\label{sec:parabolic-fermi-compare}

\begin{modelbox}[题目：两块纯金属的色散关系与输运性质比较]
设有两块纯金属，价电子能带色散关系分别为
\begin{equation}
\text{金属 I:}\quad E(k)=Ak^2,\qquad
\text{金属 II:}\quad E(k)=Bk^2,
\end{equation}
已知 $A>B$，且它们具有相同的费米能级 $E_F$。

(1) 试问它们的费米速度哪个大？
(2) 在弛豫时间相等的情况下，试问哪块金属的电导率比较大？
\end{modelbox}

\begin{tipbox}[考点快速分析]
\begin{itemize}
    \item \textbf{费米速度}：$v_F=\hbar^{-1}(dE/dk)|_{k_F}$
    \item \textbf{有效质量}：$m^*=\hbar^2(d^2E/dk^2)^{-1}$
    \item \textbf{电子浓度}：$n\propto k_F^3\propto E_F^{3/2}/A^{3/2}$
    \item \textbf{电导率}：$\sigma=ne^2\tau/m^*\propto n/m^*$
\end{itemize}
\end{tipbox}

\begin{derivationbox}[标准解题过程]
\textbf{(1) 费米速度的比较}

由抛物线色散 $E=Ak^2$，费米波矢由 $E_F$ 决定：
\begin{equation}
k_F=\sqrt{\frac{E_F}{A}}.
\end{equation}

费米速度
\begin{equation}
v_F=\frac{1}{\hbar}\left|\frac{dE}{dk}\right|_{k_F}=\frac{2Ak_F}{\hbar}=\frac{2A}{\hbar}\sqrt{\frac{E_F}{A}}=\frac{2}{\hbar}\sqrt{AE_F}.
\end{equation}

因此两金属的费米速度分别为
\begin{equation}
v_F^{\text{I}}=\frac{2}{\hbar}\sqrt{AE_F},\qquad
v_F^{\text{II}}=\frac{2}{\hbar}\sqrt{BE_F}.
\end{equation}

因 $A>B$，故
\begin{equation}
\boxed{v_F^{\text{I}}>v_F^{\text{II}}.}
\end{equation}

\textbf{(2) 电导率的比较}

有效质量
\begin{equation}
m^*=\hbar^2\left(\frac{d^2E}{dk^2}\right)^{-1}=\frac{\hbar^2}{2A},
\end{equation}
故
\begin{equation}
m_I^*=\frac{\hbar^2}{2A},\qquad m_{II}^*=\frac{\hbar^2}{2B},
\qquad\Rightarrow\qquad m_I^*<m_{II}^*.
\end{equation}

电子数密度（考虑自旋简并）
\begin{equation}
n=\frac{2}{(2\pi)^3}\cdot\frac{4\pi}{3}k_F^3=\frac{k_F^3}{3\pi^2}
=\frac{1}{3\pi^2}\left(\frac{E_F}{A}\right)^{3/2},
\end{equation}
故
\begin{equation}
n_I=\frac{1}{3\pi^2}\left(\frac{E_F}{A}\right)^{3/2},\qquad
n_{II}=\frac{1}{3\pi^2}\left(\frac{E_F}{B}\right)^{3/2},
\qquad\Rightarrow\qquad n_I<n_{II}.
\end{equation}

电导率 $\sigma=ne^2\tau/m^*$。将 $n$、$m^*$ 代入：
\begin{align}
\sigma_I&=\frac{n_I e^2\tau}{m_I^*}
=\frac{e^2\tau}{3\pi^2}\left(\frac{E_F}{A}\right)^{3/2}\cdot\frac{2A}{\hbar^2}
=\frac{2e^2\tau E_F^{3/2}}{3\pi^2\hbar^2}\cdot\frac{1}{\sqrt{A}},\\
\sigma_{II}&=\frac{2e^2\tau E_F^{3/2}}{3\pi^2\hbar^2}\cdot\frac{1}{\sqrt{B}}.
\end{align}

因 $A>B$，有 $1/\sqrt{A}<1/\sqrt{B}$，故
\begin{equation}
\boxed{\sigma_I<\sigma_{II}.}
\end{equation}
\end{derivationbox}

\begin{formulabox}[答案汇总]
\begin{center}
\begin{tabular}{lll}
\toprule
\textbf{问题} & \textbf{结论} & \textbf{原因} \\
\midrule
(1) 费米速度 & $v_F^{\text{I}}>v_F^{\text{II}}$ & 色散更陡，$v_F\propto\sqrt{A}$ \\
(2) 电导率 & $\sigma_I<\sigma_{II}$ & $n/m^*\propto A^{-1/2}$，能带窄者浓度更高 \\
\bottomrule
\end{tabular}
\end{center}
\end{formulabox}

\begin{insightbox}[物理图像]
\begin{itemize}
    \item 色散系数 $A$ 越大，能带越宽，有效质量越小，费米速度越大。但与此同时，在相同费米能下，费米球更小，电子浓度更低。
    \item 电导率正比于 $n/m^*$。对于抛物线能带，$n\propto A^{-3/2}$，$m^*\propto A^{-1}$，因此 $n/m^*\propto A^{-1/2}$。色散系数越小（能带越窄），电导率反而越大——因为电子浓度的提升超过了有效质量增加的不利影响。
    \item 这说明有效质量小并不总是意味着高电导率，电子浓度在输运性质中起着同等重要的作用。
\end{itemize}
\end{insightbox}
